\section{Reviewer-Robustness Experiments}
\label{app:reviewer-robustness}

All studies in this Online Resource section are post-hoc robustness or sensitivity analyses.
Their task selections, seeds, factors, and output paths were frozen before
formal execution, and no result changes the original 56-task discovery set,
84-task held-out set, $\alpha=0.10$ selection, or preregistered H1/H2 family.
An audit confirmed that the original Phase-2 inference already aggregates
tasks within base graph, bootstraps whole graph units, applies graph-level
sign flips, and uses the frozen Holm correction. It therefore requires no
pseudo-replication correction. New intervals below likewise aggregate paired
task effects within graph and resample graphs with equal weight.

The complete protocol grid is listed in \cref{tab:reviewer-protocols}.
% Generated by paper_scripts/reviewer_robustness/build_revision_assets.py.
\begin{table}[t]
\centering
\caption{Frozen post-hoc reviewer-robustness protocols. All task selections
used graph and structural variables only. Additional tail fractions and all
robustness intervals are exploratory and do not enter the preregistered H1/H2
family.}
\label{tab:reviewer-protocols}
\small
\begin{tabularx}{\textwidth}{@{}l>{\raggedright\arraybackslash}p{2.6cm}>{\raggedright\arraybackslash}X@{}}
\toprule
Study & Scope & Frozen factors \\
\midrule
Optimizer & 56 discovery tasks, 10 graphs & COBYLA, Nelder--Mead, SLSQP; O0/O3; $p=2,3$; three seeds; 120 actual objective calls \\
Tail fraction & 56 discovery tasks, 10 graphs & $\alpha\in\{.02,.05,.10,.25,.50,1\}$; COBYLA; $p=3$; three seeds; 120 calls \\
Depth--budget & 24 discovery tasks, 10 graphs & $p\in\{2,3,4\}$; 120/240/480-call prefixes; O0/O3; COBYLA; three seeds \\
Finite-shot training & 10 held-out tasks, 10 graphs & O0/O3; $p=3$; 1,000/10,000 shots per call; five seeds; 240-call cap \\
Classical context & All 140 canonical tasks, 25 graphs & Exact instrumented label-setting; repeated descriptive timing \\
\bottomrule
\end{tabularx}
\end{table}


\FloatBarrier
\subsection{Optimizer Robustness}

The formal matrix contains 2,016 runs: 56 tasks, three seeds, two depths, two
objectives, and three local optimizers. Each run has a strict cap of 120
actual objective calls. Wrapper and SciPy counts agree on every one of the
1,512 runs for which SciPy returns an \texttt{nfev}; the 504 SLSQP runs stop
at the wrapper cap while numerical finite differences are being evaluated, so
the absent SciPy terminal count is retained as missing rather than imputed.
Termination status, message, iteration count, wall time, and both counters are
stored per run.

Nested failures occur in both O0 and O3 under every optimizer
(\cref{tab:reviewer-optimizer}). The dual-PASS restriction then retains only
task--seed pairs for which neither arm has a positive signed nested regret
above $10^{-10}$. The positive graph intervals and frequent O0
lower-energy/lower-feasibility pairs persist in that restricted population.
The post-hoc COBYLA/O0 count of 28/168 is a separately frozen strict-cap
rerun, not a redefinition of the historical 29/168 observation, which remains
unchanged.
This supports empirical separation of optimizer inadequacy and objective
misalignment under the diagnostic, not a universal statement about either
optimizer or loss landscape.

% Generated by paper_scripts/reviewer_robustness/build_revision_assets.py.
\begin{table}[t]
\centering
\caption{Optimizer robustness on 56 discovery tasks and three fixed seeds.
Nested failures compare the best $p=3$ point with the exactly embedded $p=2$
point under 120 actual objective calls. The last two columns retain only pairs
for which both O0 and O3 pass that diagnostic; the interval is a 95\% graph-
cluster bootstrap CI for O3-minus-O0 $\gfeas$. Configurations are diagnostic,
not an optimizer ranking.}
\label{tab:reviewer-optimizer}
\small
\setlength{\tabcolsep}{4pt}
\begin{tabular}{@{}lrrrrl@{}}
\toprule
Optimizer & \shortstack{O0 nested\\failure} & \shortstack{O3 nested\\failure}
& \shortstack{dual-PASS\\pairs} & \shortstack{O0 lower energy,\\lower $\pfeas$}
& \shortstack{PASS-only $\Delta\gfeas$\\graph CI} \\
\midrule
COBYLA & 28/168 & 38/168 & 113 & 103/113 & $[0.5770,0.9971]$ \\
Nelder--Mead & 54/168 & 61/168 & 72 & 65/72 & $[0.5796,1.2473]$ \\
SLSQP & 46/168 & 72/168 & 70 & 58/70 & $[0.6096,1.3458]$ \\
\bottomrule
\end{tabular}
\end{table}


\Cref{fig:reviewer-optimizer} visualizes the failure rates and PASS-only
effects.
\begin{figure}[t]
  \centering
  \includegraphics[width=\textwidth]{fig13_reviewer_optimizer_robustness.pdf}
  \caption{Optimizer robustness. Nested-failure rates are shown separately
  for O0 and O3; paired graph-level O3-minus-O0 feasibility effects are
  evaluated on runs passing both diagnostics. All configurations use the same
  120-call resource unit, but the figure is not an optimizer leaderboard}
  \label{fig:reviewer-optimizer}
\end{figure}

\FloatBarrier
\subsection{CVaR Tail-Fraction Sensitivity}

The discovery scan completes 1,008/1,008 planned runs with no failed records.
The response peaks at the historically frozen $\alpha=0.10$ in the
equal-weight graph mean, but the nearby 0.05 and 0.25 responses differ by only
$-0.0386$ and $-0.0329$ decades; both graph-cluster intervals for these
differences include zero. By contrast, the optimized response is not assumed
to interpolate monotonically, and 54/56 task trajectories are non-monotone.
Additional held-out scans at $\alpha=0.05,0.25,0.50,1.00$ are labelled
\texttt{POST\_HOC\_SENSITIVITY}; none is used to select a new primary tail
fraction.

The $\alpha=1$ control covers 222 executed discovery and post-hoc held-out
cells and gives a maximum direct objective discrepancy of
$2.42\times10^{-13}$ from mean energy. For the 168 discovery cells matched to
the closest historical O0 control, task, depth, seed, exact embedded-$p=2$
initial parameters, nominal budget, optimizer family, initial trust-region
radius, and constraint tolerance agree. The initial objective discrepancy is
at most $2.02\times10^{-13}$. The historical and post-hoc wrappers and SciPy
versions differ, however, and complete comparable trajectories were not
persisted; their terminal parameters need not coincide. This audit therefore
validates objective equivalence only.

\Cref{tab:reviewer-alpha} reports the numerical tail-fraction summary.
% Generated by paper_scripts/reviewer_robustness/build_revision_assets.py.
\begin{table}[t]
\centering
\caption{Post-hoc CVaR tail-fraction sensitivity on 56 discovery tasks.
Values are equal-weight means over ten graph aggregates. The frozen
confirmatory choice remains $\alpha=0.10$; the final column is descriptive.}
\label{tab:reviewer-alpha}
\small
\begin{tabular}{@{}rrr@{}}
\toprule
$\alpha$ & Graph-mean $\gfeas$ & Difference from $\alpha=0.10$ \\
\midrule
0.02 & 1.8982 & -0.0907 \\
0.05 & 1.9503 & -0.0386 \\
0.10 & 1.9889 & +0.0000 \\
0.25 & 1.9560 & -0.0329 \\
0.50 & 1.8445 & -0.1444 \\
1.00 & 1.6867 & -0.3022 \\
\bottomrule
\end{tabular}
\end{table}


\begin{figure}[t]
  \centering
  \includegraphics[width=\textwidth]{fig14_reviewer_alpha_sensitivity.pdf}
  \caption{Post-hoc tail-fraction sensitivity. Discovery task summaries and
  equal-weight graph effects are displayed across the frozen grid. The curve
  is allowed to be non-monotone; $\alpha=0.10$ remains the sole confirmatory
  choice used in the original held-out family}
  \label{fig:reviewer-alpha}
\end{figure}

\FloatBarrier
\subsection{Depth--Budget Ablation}

The 24 tasks were selected using graph, size, feasible fraction, and dilution,
without reference to O0/O3 outcomes. The formal run stores one deterministic
480-call trajectory only after a smoke gate confirmed that its first 120 and
240 objective calls match separately capped runs exactly. Each reported
checkpoint is the best evaluated incumbent in that genuine prefix, rather
than necessarily the solver's last evaluated point.

% Generated by paper_scripts/reviewer_robustness/build_revision_assets.py.
\begin{table}[t]
\centering
\caption{Depth--budget ablation on the pre-frozen 24-task subset. Failure
counts pool O0 and O3 over three seeds. $p=4-p=3$ effects are equal-weight
graph means in $\gfeas$; the last column gives positive-graph counts for
O0; O3. Checkpoints are the best evaluated incumbents in verified trajectory
prefixes.}
\label{tab:reviewer-depth-budget}
\small
\begin{tabular}{@{}rrrrrr@{}}
\toprule
$N_{\rm fev}$ & $p2\to p3$ failure & $p3\to p4$ failure & \shortstack{$p4-p3$\\O0 $\gfeas$} & \shortstack{$p4-p3$\\O3 $\gfeas$} & \shortstack{positive graphs\\O0; O3} \\
\midrule
120 & 30/144 (20.8\%) & 45/144 (31.2\%) & 0.7206 & 0.2021 & 10/10; 10/10 \\
240 & 31/144 (21.5\%) & 45/144 (31.2\%) & 0.7843 & 0.2074 & 10/10; 10/10 \\
480 & 38/144 (26.4\%) & 47/144 (32.6\%) & 0.8683 & 0.1970 & 10/10; 10/10 \\
\bottomrule
\end{tabular}
\end{table}


The failure-rate direction and aggregation were independently audited. All 16
structural checks passed, and 24 stratified PASS/FAIL cases were recomputed
from their saved parameters. Maximum absolute discrepancies were
$1.55\times10^{-15}$ for the deeper objective,
$2.22\times10^{-16}$ for the embedded objective, and
$1.42\times10^{-15}$ for signed regret; both embedding identities were exact
at stored precision. The licensed conclusion is limited to the tested subset,
COBYLA, O0/O3, depths, seeds, and budgets: increasing the tested classical
evaluation budget did not monotonically reduce certified nested-ansatz
failures. The positive $p=4-p=3$ feasibility effects in
\cref{tab:reviewer-depth-budget} show why that
diagnostic must not be conflated with final-state quality.

\FloatBarrier
\subsection{Finite-Shot Estimation}

The all-held-out fixed-endpoint protocol and its 20 replicates per cell are
detailed in \cref{app:finite-shot}. A separate outcome-blind 18-task subset
uses 50 sampling seeds to cross
$\alpha\in\{0.02,0.05,0.10,0.25,0.50\}$ with
$N_{\rm shots}\in\{10^3,10^4,10^5\}$. The mean of the separately computed
O0 and O3 CVaR-0.10 RMSE values falls from 0.04839 to 0.01570 and 0.00486 as
shots increase. The heat map in \cref{fig:reviewer-alpha-shots} shows the larger
sampling cost of more aggressive lower tails. No quantum optimization is
rerun for this estimator analysis.

\begin{figure}[t]
  \centering
  \includegraphics[width=0.92\textwidth]{fig16_reviewer_alpha_shots.pdf}
  \caption{CVaR estimator robustness across tail fraction and shot count on
  fixed exact-trained endpoints. Cells report sampling error, not retrained
  circuit quality. Smaller tails generally require more shots for comparable
  precision}
  \label{fig:reviewer-alpha-shots}
\end{figure}

\FloatBarrier
\subsection{Finite-Shot Training}
\label{app:reviewer-finite-shot}

The finite-shot training matrix completes 200/200 runs with no failed records.
Each objective evaluation uses a fresh reproducible sample; final parameters
are evaluated exactly. The two shot-trained point estimates preserve the
exact O3-minus-O0 direction, but both graph-cluster intervals cross zero
(\cref{tab:reviewer-finite-training}). The 10,000-shot condition has positive
effects on eight of ten graphs, compared with six of ten at 1,000 shots and
nine of ten for matched exact training. This is mixed evidence about noisy
training and stronger evidence about fixed-endpoint estimation; it is not a
hardware or device-noise experiment.

% Generated by paper_scripts/reviewer_robustness/build_revision_assets.py.
\begin{table}[t]
\centering
\caption{O3-minus-O0 finite-shot training effects on ten pre-frozen tasks
from ten held-out graphs. Final parameters are evaluated exactly; intervals
are 95\% graph-cluster bootstrap CIs. The exact row uses the matched canonical
endpoints for the same tasks.}
\label{tab:reviewer-finite-training}
\small
\begin{tabular}{@{}lrrl@{}}
\toprule
Training regime & Positive graphs & Mean $\Delta\gfeas$ & Graph CI \\
\midrule
Exact training & 9/10 & 0.3599 & $[0.1641,0.5697]$ \\
10,000-shot training & 8/10 & 0.1706 & $[-0.0106,0.3862]$ \\
1,000-shot training & 6/10 & 0.1246 & $[-0.0050,0.3351]$ \\
\bottomrule
\end{tabular}
\end{table}


The main article's robustness figure displays these graph intervals and
terminal-state distributions; the complete numerical matrix remains in
\cref{tab:reviewer-finite-training}.

\FloatBarrier
\subsection{Classical RCSP Context}

An instrumented exact label-setting solver reproduces the frozen optimum on
140/140 canonical tasks. Across repeated high-resolution solves, the median,
95th-percentile, and maximum per-task medians are 50.0, 95.4, and
188.1 microseconds; generated-label counts have median 13 and maximum 30.
These small descriptive values confirm the intended context: the instances
are classically tractable mechanism probes. They do not define a
quantum--classical runtime comparison.

\Cref{tab:reviewer-classical} gives the compact classical-solver summary.
% Generated by paper_scripts/reviewer_robustness/build_revision_assets.py.
\begin{table}[t]
\centering
\caption{Exact classical RCSP context. Times are descriptive medians of
repeated high-resolution measurements and are not compared with statevector
runtime or hardware performance.}
\label{tab:reviewer-classical}
\scriptsize
\setlength{\tabcolsep}{4pt}
\begin{tabular}{@{}rrrrrr@{}}
\toprule
Matches & Median ($\mu$s) & p95 ($\mu$s) & Max ($\mu$s) & \shortstack{Median\\labels} & \shortstack{Maximum\\labels} \\
\midrule
140/140 & 50.0 & 95.4 & 188.1 & 13 & 30 \\
\bottomrule
\end{tabular}
\end{table}


\FloatBarrier
